\chapter{Use Cases}
\ChapterLabel{UseCases}
\label{Section:UseCasesResponsibility}
This chapter describes live measurement delivery and diagnostic recording and playback through six candidate scenarios, with actors, flows, requirements, planned evidence and open decisions.

\section{Scenario Scope}
The use-case boundary is the sensor processing subsystem. The Sensor/FTDI
hardware and the Mission application are external to this boundary.
The two live measurement scenarios send output to the Mission application using the Open Mission
Systems (OMS) interface selected for this project. Record the applicable
OMS baseline and the sample and status contracts in
\hyperref[Question:Q-114]{Q-114}.


Diagnostic scenarios retain the sensor processing subsystem boundary and add the engineering actors, recording storage and selected replay consumer. \gls{diagnosticrecording} preserves the observations available at the selected capture boundary; \gls{playback} supplies those recorded inputs for investigation. These scenarios and diagrams are proposed behavior, not selected software components or an executed product. No new hardware platform or GUI is selected.
\index{diagnostic recording}\index{playback}
Figure~\ref{Figure:DiagnosticUseCases} introduces \UseCaseRef{UC-03}, \UseCaseRef{UC-04}, \UseCaseRef{UC-05}, and \UseCaseRef{UC-06}.
\begin{figure}[htbp]\centering
\resizebox{0.98\linewidth}{!}{\begin{tikzpicture}[>=Latex,font=\small,
 actor/.style={draw,rounded corners,align=center,text width=2.8cm,minimum height=0.9cm},
 goal/.style={draw,ellipse,align=center,text width=3.6cm,minimum height=1.25cm}]
\node[actor] (field) at (0,0) {\hyperref[Stakeholder:Field-Engineer]{Field-Engineer}};
\node[actor] (devopen) at (0,-3) {\hyperref[Stakeholder:Acquisition-Developer]{Acquisition-Developer}};
\node[actor] (devplay) at (0,-6) {\hyperref[Stakeholder:Acquisition-Developer]{Acquisition-Developer}};
\node[actor] (inspect) at (0,-9) {\hyperref[Stakeholder:Field-Engineer]{Field-Engineer}};
\node[goal] (record) at (5.4,0) {\hyperref[UseCase:UC-03]{UC-03\\Record a Diagnostic Session}};
\node[goal] (open) at (5.4,-3) {\hyperref[UseCase:UC-04]{UC-04\\Open and Validate a Recording}};
\node[goal] (play) at (5.4,-6) {\hyperref[UseCase:UC-05]{UC-05\\Replay a Diagnostic Session}};
\node[goal] (control) at (5.4,-9) {\hyperref[UseCase:UC-06]{UC-06\\Inspect Events During Playback}};
\node[actor] (integrator) at (10.7,-6) {\hyperref[Stakeholder:Mission-Integrator]{Mission-Integrator}};
\node[draw,dashed,fit=(record)(open)(play)(control),inner sep=0.32cm,
 label=above:{Sensor processing subsystem}] {};
\draw (field.east) -- (record.west);
\draw (devopen.east) -- (open.west);
\draw (devplay.east) -- (play.west);
\draw (inspect.east) -- (control.west);
\draw (play.east) -- (integrator.west);
\node[align=center] at (5.4,-11.2) {Repeated actor labels denote the same stakeholder role.};
\end{tikzpicture}
}
\caption{Engineering actors and recording/playback goals at the subsystem boundary.}\label{Figure:DiagnosticUseCases}
\end{figure}

\section{Receive Qualified Sensor Samples}

\small
\begin{tabularx}{\linewidth}{L{2.8cm}Y}
\caption{Receive Qualified Sensor Samples\label{Table:UC-01}\label{Table:SensorUseCase}}\\
\toprule
Field & Candidate scenario\\\midrule
\endfirsthead
\toprule
Field & Candidate scenario\\\midrule
\endhead
\bottomrule
\endlastfoot
ID & \UseCaseID{UC-01}\\
External input & Sensor/FTDI hardware delivers a conversion code and its associated counter value.\\
External recipient & Mission application via OMS.\\
Source stories & \UserStoryRef{US-01}, \UserStoryRef{US-02}, \UserStoryRef{US-04}.\\
Requirements & \RequirementRef{REQ-F-101}, \RequirementRef{REQ-F-102}, \RequirementRef{REQ-F-103}, \RequirementRef{REQ-C-101}, \RequirementRef{REQ-C-102}, \RequirementRef{REQ-C-103}, \RequirementRef{REQ-F-104}, \RequirementRef{REQ-F-105}, \RequirementRef{REQ-F-106}, \RequirementRef{REQ-F-107}, \RequirementRef{REQ-I-101}, \RequirementRef{REQ-SMP-01}, \RequirementRef{REQ-TIM-01}, \RequirementRef{REQ-IF-01}, \RequirementRef{REQ-TIM-02}, \RequirementRef{REQ-F-110}\\
Tests & \TestCaseRef{TC-101}, \TestCaseRef{TC-102}, \TestCaseRef{TC-103}, \TestCaseRef{TC-104}, \TestCaseRef{TC-105}, \TestCaseRef{TC-106}, \TestCaseRef{TC-107}, \TestCaseRef{TC-108}, \TestCaseRef{TC-109}, \TestCaseRef{TC-110}, \TestCaseRef{TC-113}, \TestCaseRef{TC-114}, \TestCaseRef{TC-115}, \TestCaseRef{TC-116}, \TestCaseRef{TC-118}, \TestCaseRef{TC-120}\\
Preconditions & The hardware link, capture rules, qualification rules, and OMS sample contract are configured for the scenario.\\
Trigger & A sample arrives from the Sensor/FTDI hardware.\\
Main flow & The subsystem receives the sample and associated counter value, determines measurement usability, and exposes the local sample record under \RequirementRef{REQ-IF-01}. The Mission application receives the qualified value, status, and timing over OMS. The producer and contract for mission sample delivery are open in \hyperref[Question:Q-118]{Q-118}. The Operator sees the reading attributes at the agreed point of use under \RequirementRef{REQ-F-110}. As earlier samples are processed, acquisition continues as called for by \UserStoryRef{US-02}.\\
Success outcome & The Mission application receives and interprets the qualified sample according to the agreed OMS sample contract.\\
Alternate flow & If a reading is unusable, apply the qualification and publication rules agreed in \hyperref[Question:Q-113]{Q-113} and \hyperref[Question:Q-114]{Q-114}. The disposition remains open.\\
Exceptions & Partial records, loss of input, and delivery failures require agreed handling in \hyperref[Question:Q-107]{Q-107} and \hyperref[Question:Q-114]{Q-114}. Stream cadence and allowable gaps remain open in \hyperref[Question:Q-104]{Q-104}.\\
Postconditions & The Mission application has the sample and its qualification information for the successful transfer.\\
\end{tabularx}
\normalsize

\subsection{Quality Scenarios}
These candidate measures identify what to observe. The linked questions own
the operating limits, expected outputs, and pass thresholds.

\subsubsection{QS-01 Sample integrity and qualification}
\label{Scenario:QS-01}
\begin{description}[leftmargin=0pt,labelindent=0pt,labelsep=0pt,style=nextline]
\item[Stimulus source] Sensor/FTDI hardware driven by reference inputs.
\item[Stimulus] Delivers known conversion-code and counter pairs, including qualification boundary cases.
\item[Environment] Configured acquisition link under nominal operation and qualification boundary conditions.
\item[Artifact] Capture, sample association, qualification, and OMS sample output.
\item[Response] Preserve the sample-to-counter association and publish the value and qualification disposition specified for each reference case.
\item[Response measure] Use the reference vectors to check received values, counter associations, and qualification dispositions. Count mismatches by field and case. Reference inputs and limits are open in \hyperref[Question:Q-108]{Q-108}, \hyperref[Question:Q-112]{Q-112}, and \hyperref[Question:Q-113]{Q-113}.
\end{description}

\subsubsection{QS-02 Acquisition continuity under processing load}
\label{Scenario:QS-02}
\begin{description}[leftmargin=0pt,labelindent=0pt,labelsep=0pt,style=nextline]
\item[Stimulus source] Sensor/FTDI hardware.
\item[Stimulus] Delivers new samples while earlier samples are still being processed.
\item[Environment] Sustained acquisition under the supported processing load.
\item[Artifact] Acquisition, processing, and OMS sample delivery path.
\item[Response] Continue acquiring new samples and deliver the qualified stream to the Mission application.
\item[Response measure] Measure delivered sample rate, maximum gap between deliveries, and count of missing expected samples over the agreed observation interval. Check the 10 ms local publication allocation in \RequirementRef{REQ-TIM-01}. Stream load and delivery limits are open in \hyperref[Question:Q-104]{Q-104}.
\end{description}

\subsubsection{QS-03 Sample freshness}
\label{Scenario:QS-03}
\begin{description}[leftmargin=0pt,labelindent=0pt,labelsep=0pt,style=nextline]
\item[Stimulus source] Sensor/FTDI hardware and the sample processing path.
\item[Stimulus] A sample reaches qualification or delivery after a delay near or beyond the freshness limit.
\item[Environment] Active acquisition with representative processing and delivery delay.
\item[Artifact] Sample timing, qualification, and OMS sample output.
\item[Response] Preserve timing information and apply the agreed freshness disposition so the Mission application can interpret sample age and usability.
\item[Response measure] Measure capture-to-receipt latency and reported or derived age error against reference timestamps. Compare dispositions at the freshness boundary with the agreed examples and check local event timestamp error against the 1 ms allocation in \RequirementRef{REQ-TIM-02}. Timing origin and mission delivery limits are open in \hyperref[Question:Q-104]{Q-104}, \hyperref[Question:Q-113]{Q-113}, and \hyperref[Question:Q-114]{Q-114}.
\end{description}

\subsubsection{QS-04 Sample stream degradation and recovery}
\label{Scenario:QS-04}
\begin{description}[leftmargin=0pt,labelindent=0pt,labelsep=0pt,style=nextline]
\item[Stimulus source] Sensor/FTDI hardware or the OMS delivery path.
\item[Stimulus] Input becomes partial or unavailable, or delivery is interrupted.
\item[Environment] An established sample stream during a fault and subsequent restoration.
\item[Artifact] Acquisition, qualification, and OMS sample delivery path.
\item[Response] Apply the agreed disposition for affected readings and the agreed policy for resuming delivery.
\item[Response measure] Compare affected readings and observable outputs with expected fault outcomes. Measure time from fault onset to observable disposition and from restoration to the first accepted sample, where applicable. Outcomes and limits are open in \hyperref[Question:Q-107]{Q-107} and \hyperref[Question:Q-114]{Q-114}.
\end{description}

\subsubsection{QS-05 Sample interface substitutability}
\label{Scenario:QS-05}
\begin{description}[leftmargin=0pt,labelindent=0pt,labelsep=0pt,style=nextline]
\item[Stimulus source] Mission-Integrator.
\item[Stimulus] Independently replaces the sensor or Mission application with a compatible implementation.
\item[Environment] Integration using the documented sensor-side and OMS sample contracts.
\item[Artifact] Sensor-side sample interface and OMS sample output.
\item[Response] Accept equivalent sensor input or deliver samples the replacement Mission application can interpret using the existing contracts.
\item[Response measure] Run the same sample contract cases before and after each swap. Record field or interpretation mismatches and whether either contract or connection design needed to change. Compatibility criteria are open in \hyperref[Question:Q-102]{Q-102} and \hyperref[Question:Q-114]{Q-114}.
\end{description}

\section{Monitor Sensor Status}

\small
\begin{tabularx}{\linewidth}{L{2.8cm}Y}
\caption{Monitor Sensor Status\label{Table:UC-02}}\\
\toprule
Field & Candidate scenario\\\midrule
\endfirsthead
\toprule
Field & Candidate scenario\\\midrule
\endhead
\bottomrule
\endlastfoot
ID & \UseCaseID{UC-02}\\
Status source & The subsystem uses observed source connection and operating-health information.\\
External recipient & Mission application via OMS.\\
Source stories & \UserStoryRef{US-01} for availability and \UserStoryRef{US-04} for the documented status interface.\\
Requirements & \RequirementRef{REQ-F-108}, \RequirementRef{REQ-F-109}, \RequirementRef{REQ-SMP-02}, \RequirementRef{REQ-OMS-01}\\
Tests & \TestCaseRef{TC-111}, \TestCaseRef{TC-112}, \TestCaseRef{TC-117}, \TestCaseRef{TC-119}\\
Preconditions & Source-health criteria and the OMS status contract are defined for the scenario.\\
Trigger & A source-status observation becomes available for reporting. Reporting cadence and change triggers remain open in \hyperref[Question:Q-114]{Q-114}.\\
Main flow & The subsystem determines source connection and operating-health status from available observations and reports derived status to the Mission application. The mapping from observations to valid, stale, or faulted status remains open in \hyperref[Question:Q-125]{Q-125}.\\
Success outcome & The Mission application can interpret the reported source condition using the agreed OMS status contract.\\
Alternate flow & When observations are stale or insufficient, apply the agreed reporting policy. Status vocabulary and handling remain open in \hyperref[Question:Q-115]{Q-115}.\\
Exceptions & Missing observations and loss of delivery require agreed handling in \hyperref[Question:Q-107]{Q-107} and \hyperref[Question:Q-114]{Q-114}.\\
Postconditions & Source status is available to the Mission application for the successful transfer. The acquisition service exposes stale state and the mission gateway reports derived status under the candidate contracts in \RequirementRef{REQ-SMP-02} and \RequirementRef{REQ-OMS-01}.\\
\end{tabularx}
\normalsize

\subsection{Quality Scenarios}

\subsubsection{QS-06 Status accuracy and delivery timeliness}
\label{Scenario:QS-06}
\begin{description}[leftmargin=0pt,labelindent=0pt,labelsep=0pt,style=nextline]
\item[Stimulus source] Sensor source and connection observations.
\item[Stimulus] A reference fixture produces a known connection transition or operating-health change.
\item[Environment] Monitoring is active and OMS delivery is available.
\item[Artifact] Source-status determination and OMS status output.
\item[Response] Determine the source condition and deliver the corresponding status to the Mission application according to the reporting policy.
\item[Response measure] Compare each delivered state against independent fixture evidence. Measure time from the effective source transition to receipt by the Mission application. Compare CAL messages for valid, stale, and faulted states against the mapping in \RequirementRef{REQ-OMS-01} after program approval. State criteria, reporting cadence, and latency limits are open in \hyperref[Question:Q-114]{Q-114} and \hyperref[Question:Q-115]{Q-115}.
\end{description}

\subsubsection{QS-07 Stale or insufficient status evidence}
\label{Scenario:QS-07}
\begin{description}[leftmargin=0pt,labelindent=0pt,labelsep=0pt,style=nextline]
\item[Stimulus source] Source-status observations.
\item[Stimulus] Observations age past the agreed limit, disappear, or become insufficient to determine source condition.
\item[Environment] Monitoring continues after a previously confirmed state.
\item[Artifact] Status freshness evaluation and OMS status output.
\item[Response] Apply the agreed disposition for stale, missing, or insufficient evidence.
\item[Response measure] Compare the delivered status with the expected fault disposition. Measure elapsed time from the evidence timeout or loss to the observable disposition where applicable. Check the 100 ms local stale-state allocation in \RequirementRef{REQ-SMP-02} separately from OMS delivery. Evidence criteria and reporting policy are open in \hyperref[Question:Q-107]{Q-107}, \hyperref[Question:Q-114]{Q-114}, and \hyperref[Question:Q-115]{Q-115}.
\end{description}

\subsubsection{QS-08 Status delivery interruption and recovery}
\label{Scenario:QS-08}
\begin{description}[leftmargin=0pt,labelindent=0pt,labelsep=0pt,style=nextline]
\item[Stimulus source] OMS delivery path.
\item[Stimulus] Status delivery is interrupted and later becomes available.
\item[Environment] Monitoring remains active while source observations continue.
\item[Artifact] OMS status publication path and Mission application receipt.
\item[Response] Apply the agreed reporting and recovery policy during and after the interruption.
\item[Response measure] Measure the reporting gap and the time from delivery restoration to receipt of the first status allowed by the recovery policy. Compare its state and age with reference observations. Recovery policy and limits are open in \hyperref[Question:Q-107]{Q-107} and \hyperref[Question:Q-114]{Q-114}.
\end{description}

\subsubsection{QS-09 Status interface substitutability}
\label{Scenario:QS-09}
\begin{description}[leftmargin=0pt,labelindent=0pt,labelsep=0pt,style=nextline]
\item[Stimulus source] Mission-Integrator.
\item[Stimulus] Replaces the Mission application with a compatible OMS status consumer.
\item[Environment] Integration using the documented OMS status contract and representative source states.
\item[Artifact] OMS status output and Mission application interpretation.
\item[Response] Deliver status the replacement Mission application can interpret using the existing contract.
\item[Response measure] Run the same status contract cases before and after the swap. Record state, freshness, and interpretation mismatches and whether the contract or connection design needed to change. Compatibility criteria are open in \hyperref[Question:Q-114]{Q-114}.
\end{description}

\begin{figure}[htbp]\centering
\begin{tikzpicture}[>=Latex]
\node[draw,align=center,text width=2.6cm] (sensor) at (0,1.3)
  {Sensor/FTDI hardware};
\node[draw,ellipse,align=center,text width=3.3cm] (sample) at (4.8,1.3)
  {\hyperref[UseCase:UC-01]{UC-01\\Receive Qualified Sensor Samples}};
\node[draw,ellipse,align=center,text width=3.3cm] (status) at (4.8,-1.3)
  {\hyperref[UseCase:UC-02]{UC-02\\Monitor Sensor Status}};
\node[draw,dashed,fit=(sample)(status),inner sep=0.35cm,
  label=above:{Sensor processing subsystem}] {};
\node[draw,align=center,text width=2.3cm,minimum height=3.2cm] (mission) at (10,0)
  {Mission application};
\draw[->] (sensor.east) -- (sample.west);
\draw[->] (sample.east) -- node[above,pos=0.62]{OMS} (mission.west |- sample.east);
\draw[->] (status.east) -- node[above,pos=0.62]{OMS} (mission.west |- status.east);
\end{tikzpicture}

\caption{Sensor/FTDI input and OMS outputs at the sensor processing subsystem boundary.}\label{Figure:UseCaseDiagram}
\end{figure}

\section{Record a Diagnostic Session}
\par\medskip
\small
\begin{tabularx}{\linewidth}{L{2.8cm}Y}
\caption{Record a Diagnostic Session\label{Table:UC-03}}\\
\toprule
Field & Candidate scenario\\\midrule
\endfirsthead
\toprule
Field & Candidate scenario\\\midrule
\endhead
\bottomrule
\endlastfoot
ID & \UseCaseID{UC-03}\\
Actors & Field-Engineer; Acquisition-Developer supplies reference inputs during verification.\\
Source stories & \UserStoryRef{US-05}, \UserStoryRef{US-08}\\
Requirements & \RequirementRef{REQ-F-111}, \RequirementRef{REQ-F-112}, \RequirementRef{REQ-D-101}, \RequirementRef{REQ-D-102}, \RequirementRef{REQ-F-113}, \RequirementRef{REQ-D-103}, \RequirementRef{REQ-Q-102}, \RequirementRef{REQ-Q-103}, \RequirementRef{REQ-Q-104}, \RequirementRef{REQ-Q-105}\\
Tests & \TestCaseRef{TC-121}, \TestCaseRef{TC-122}, \TestCaseRef{TC-123}, \TestCaseRef{TC-124}, \TestCaseRef{TC-125}, \TestCaseRef{TC-126}, \TestCaseRef{TC-144}, \TestCaseRef{TC-145}, \TestCaseRef{TC-146}, \TestCaseRef{TC-147}\\
Validation & \ValidationRef{VAL-121}\\
Preconditions & The recording boundary, schema, context metadata and fault policy are defined; storage and representative input are available.\\
Trigger & The Field-Engineer requests a diagnostic recording.\\
Main flow & 1. The Field-Engineer starts recording. 2. The subsystem records the observed samples, results and diagnostic events with capture context. 3. Recording continues under the agreed live load. 4. The Field-Engineer requests stop. 5. The subsystem finalizes the recording under the stop-and-flush policy.\\
Success outcome & A readable incident recording retains the agreed capture interval, data associations and context for later investigation.\\
Alternate flow & Known input/recorder gaps are identified; invalid inputs retain the representation defined by the recording schema.\\
Exceptions & Storage is unavailable/full, writes fail, or capture is interrupted. Report failure and preserve recoverable content according to the policy; do not report successful completion for a failed recording.\\
Postconditions & The recording is finalized, or an explicit failure and incomplete-content disposition is available.\\
Open decisions & \hyperref[Question:Q-127]{Q-127}, \hyperref[Question:Q-131]{Q-131}. Other applicable contracts are linked from each requirement.\\
\end{tabularx}
\normalsize
Figure~\ref{Figure:RecordingActivity} shows the recording success and failure paths.
\begin{figure}[htbp]\centering
\resizebox{0.92\linewidth}{!}{\begin{tikzpicture}[>=Latex,font=\small,
 activity/.style={draw,rounded corners,align=center,text width=4.5cm,minimum height=0.85cm},
 choice/.style={draw,diamond,aspect=2.3,align=center,inner sep=2pt,text width=2.2cm}]
\node[activity] (start) at (0,0) {Field-Engineer requests recording};
\node[activity] (context) at (0,-1.5) {Create recording identity and retain capture context};
\node[activity] (capture) at (0,-3.2) {Persist observed samples, results and diagnostic events};
\node[choice] (fault) at (0,-5.6) {Recording failure?};
\node[choice] (stop) at (0,-8.2) {Stop requested?};
\node[activity] (final) at (0,-10.8) {Apply stop-and-flush policy; finalize readable recording};
\node[activity] (failed) at (6,-5.6) {Report failure and incomplete-content disposition};
\node[activity] (done) at (0,-12.6) {Report completed recording};
\draw[->] (start) -- (context);
\draw[->] (context) -- (capture);
\draw[->] (capture) -- (fault);
\draw[->] (fault) -- node[right]{No} (stop);
\draw[->] (fault.east) -- node[above]{Yes} (failed.west);
\draw[->] (stop) -- node[right]{Yes} (final);
\draw[->] (final) -- (done);
\draw[->] (stop.west) -- (-3.8,-8.2) -- node[left,align=right]{No:\\continue} (-3.8,-3.2) -- (capture.west);
\node[align=left,text width=4.4cm] at (6,-9.7) {Known gaps are retained under the recording policy. A finalization failure follows the failure path, not completion.};
\draw[->] (final.east) -- (3,-10.8) -- (3,-7.1) -- (6,-7.1) -- (failed.south);
\node[fill=white,inner sep=2pt] at (4.5,-7.1) {Finalize failed};
\end{tikzpicture}
}
\caption{Candidate recording activity, including explicit recording failure.}\label{Figure:RecordingActivity}
\end{figure}

\section{Open and Validate a Recording}
\par\medskip
\small
\begin{tabularx}{\linewidth}{L{2.8cm}Y}
\caption{Open and Validate a Recording\label{Table:UC-04}}\\
\toprule
Field & Candidate scenario\\\midrule
\endfirsthead
\toprule
Field & Candidate scenario\\\midrule
\endhead
\bottomrule
\endlastfoot
ID & \UseCaseID{UC-04}\\
Actors & Acquisition-Developer.\\
Source stories & \UserStoryRef{US-06}, \UserStoryRef{US-08}\\
Requirements & \RequirementRef{REQ-F-114}, \RequirementRef{REQ-F-115}, \RequirementRef{REQ-F-116}\\
Tests & \TestCaseRef{TC-127}, \TestCaseRef{TC-128}, \TestCaseRef{TC-129}\\
Validation & \ValidationRef{VAL-122}\\
Preconditions & Compatibility rules, format readers and the context presentation contract are defined; the selected recording is accessible.\\
Trigger & The Acquisition-Developer selects a recorded session.\\
Main flow & 1. Select a recording. 2. The subsystem reads its format and capture context. 3. It evaluates compatibility. 4. The Acquisition-Developer reviews the context and compatibility disposition. 5. A compatible recording becomes ready for playback.\\
Success outcome & The developer understands the capture context and knows whether the recording can be replayed.\\
Alternate flow & Missing metadata or a different software/configuration version receives the disposition specified by the compatibility rules.\\
Exceptions & Unsupported, incompatible, corrupt, truncated, or changed content is classified under the approved rules; playback is prevented when incompatible and the reason is presented.\\
Postconditions & A compatible recording is ready for playback, or an incompatibility disposition is available without admitting recorded input.\\
Open decisions & \hyperref[Question:Q-128]{Q-128}, \hyperref[Question:Q-131]{Q-131}. Other applicable contracts are linked from each requirement.\\
\end{tabularx}
\normalsize

\section{Replay a Diagnostic Session}
\par\medskip
\small
\begin{tabularx}{\linewidth}{L{2.8cm}Y}
\caption{Replay a Diagnostic Session\label{Table:UC-05}}\\
\toprule
Field & Candidate scenario\\\midrule
\endfirsthead
\toprule
Field & Candidate scenario\\\midrule
\endhead
\bottomrule
\endlastfoot
ID & \UseCaseID{UC-05}\\
Actors & Acquisition-Developer; Mission-Integrator observes the selected output consumer.\\
Source stories & \UserStoryRef{US-06}, \UserStoryRef{US-09}\\
Requirements & \RequirementRef{REQ-F-117}, \RequirementRef{REQ-F-127}, \RequirementRef{REQ-D-104}, \RequirementRef{REQ-Q-101}, \RequirementRef{REQ-I-102}, \RequirementRef{REQ-I-103}, \RequirementRef{REQ-F-126}\\
Tests & \TestCaseRef{TC-130}, \TestCaseRef{TC-149}, \TestCaseRef{TC-131}, \TestCaseRef{TC-132}, \TestCaseRef{TC-141}, \TestCaseRef{TC-142}, \TestCaseRef{TC-143}\\
Validation & \ValidationRef{VAL-123}\\
Preconditions & The recording has passed the compatibility gate; input boundary, time/order rules, replay identification and output routing contracts are defined; an allowed endpoint is selected.\\
Trigger & The Acquisition-Developer requests playback of a compatible recording.\\
Main flow & 1. Start playback without Sensor/FTDI hardware and establish the approved initial processing state before admitting the first event. 2. The subsystem admits recorded events in their defined order and timing. 3. The processing path produces results. 4. The consumer and engineer receive explicitly identified replay outputs at the selected endpoint. 5. At end of recording, completion is reported.\\
Success outcome & The developer can investigate the captured software behavior and the Mission-Integrator can distinguish recorded measurements from current live measurements.\\
Alternate flow & Recorded gaps and disconnect events are replayed as defined by the captured event schema; playback can be inspected through \UseCaseRef{UC-06}.\\
Exceptions & Unavailable initial processing state receives the approved disposition without admitting recorded input. Read, processing or delivery failure follows the defined lifecycle/routing disposition. Missing endpoint selection does not silently select a mission endpoint.\\
Postconditions & The last processed position and completion/failure state are available; outputs have followed the approved replay-identification and routing rules.\\
Open decisions & \hyperref[Question:Q-129]{Q-129}, \hyperref[Question:Q-131]{Q-131}. Other applicable contracts are linked from each requirement.\\
\end{tabularx}
\normalsize
Figure~\ref{Figure:PlaybackSequence} shows externally meaningful playback interactions. Subsystem lanes denote behavior, not selected processes, threads or classes. See \hyperref[Question:Q-129]{Q-129} and \hyperref[Question:Q-130]{Q-130} for the input and routing contracts.
\begin{figure}[htbp]\centering
\resizebox{0.98\linewidth}{!}{\begin{tikzpicture}[>=Latex,font=\small,
 lane/.style={draw,align=center,text width=2.7cm,minimum height=1.1cm}]
\node[lane] (user) at (0,0) {Acquisition-Developer};
\node[lane] (session) at (4,0) {Subsystem:\\recording access};
\node[lane] (process) at (8,0) {Subsystem:\\playback processing};
\node[lane] (consumer) at (12,0) {Selected output consumer};
\foreach \x in {0,4,8,12} {\draw[dashed] (\x,-0.7) -- (\x,-13.3);}
\draw[->] (0,-1.8) -- node[above]{Select recording} (4,-1.8);
\draw[->] (4,-3.2) -- node[above]{Context and compatibility} (0,-3.2);
\node[draw,dashed,align=center,fill=white,text width=10cm] at (6,-4.8)
 {Precondition: compatible recording and permitted endpoint selected};
\draw[->] (0,-6.4) -- node[above,pos=0.35]{Start playback; initialize state} (8,-6.4);
\draw[->] (4,-7.8) -- node[above]{Recorded events} (8,-7.8);
\node[align=center,text width=3.6cm,fill=white] at (8,-9.3) {Process in defined order and replay time};
\draw[->] (8,-11.2) -- node[above]{Identified replay outputs} (12,-11.2);
\draw[->] (8,-12.8) -- node[above,pos=0.6]{Position and completion/failure state} (0,-12.8);
\node[align=center,text width=11cm] at (6,-14.6)
 {Event delivery repeats to end of recording or failure. Sensor hardware is not required. Failure and routing alternatives are defined in the use-case table.};
\end{tikzpicture}
}
\caption{Candidate playback interaction after compatibility and endpoint selection.}\label{Figure:PlaybackSequence}
\end{figure}

\section{Inspect Events During Playback}
\par\medskip
\small
\begin{tabularx}{\linewidth}{L{2.8cm}Y}
\caption{Inspect Events During Playback\label{Table:UC-06}}\\
\toprule
Field & Candidate scenario\\\midrule
\endfirsthead
\toprule
Field & Candidate scenario\\\midrule
\endhead
\bottomrule
\endlastfoot
ID & \UseCaseID{UC-06}\\
Actors & Field-Engineer.\\
Source stories & \UserStoryRef{US-07}\\
Requirements & \RequirementRef{REQ-F-118}, \RequirementRef{REQ-F-119}, \RequirementRef{REQ-F-120}, \RequirementRef{REQ-F-121}, \RequirementRef{REQ-F-122}, \RequirementRef{REQ-F-123}, \RequirementRef{REQ-F-124}, \RequirementRef{REQ-F-125}, \RequirementRef{REQ-Q-106}\\
Tests & \TestCaseRef{TC-133}, \TestCaseRef{TC-134}, \TestCaseRef{TC-135}, \TestCaseRef{TC-136}, \TestCaseRef{TC-137}, \TestCaseRef{TC-138}, \TestCaseRef{TC-139}, \TestCaseRef{TC-140}, \TestCaseRef{TC-148}\\
Validation & \ValidationRef{VAL-124}\\
Preconditions & A compatible recording is open; supported controls, playback states, replay-clock and state-restoration policies are defined.\\
Trigger & The Field-Engineer requests a playback control action while investigating an incident.\\
Main flow & 1. Locate the incident position. 2. Pause and inspect the selected record and status. 3. Step through individual events. 4. Seek backward or forward with the defined state restoration. 5. Change speed or resume and inspect the event sequence.\\
Success outcome & The engineer can examine the events surrounding a fault with interpretable position, state and timing information.\\
Alternate flow & A seek restores the approved state before further event admission; pausing and speed changes apply the replay-clock policy to freshness and timeouts.\\
Exceptions & Invalid locations, unsupported speeds and controls issued in an invalid state receive the approved disposition. End-of-recording and processing errors follow the playback lifecycle policy.\\
Postconditions & Playback position and state correspond to the last accepted action, with restored processing state where required.\\
Open decisions & \hyperref[Question:Q-129]{Q-129}, \hyperref[Question:Q-131]{Q-131}. Other applicable contracts are linked from each requirement.\\
\end{tabularx}
\normalsize
Figure~\ref{Figure:PlaybackStates} illustrates the candidate observable states. Ready follows successful compatibility evaluation; Playing admits events, Paused permits inspection and single-step, Completed indicates end of recording, and Error indicates failed playback. Seeking and speed selection apply the policy in \hyperref[Question:Q-129]{Q-129}; the diagram does not select an internal state-machine implementation.
\begin{figure}[htbp]\centering
\resizebox{0.96\linewidth}{!}{\begin{tikzpicture}[>=Latex,font=\small,
 state/.style={draw,rounded corners,align=center,text width=2.8cm,minimum height=1cm}]
\node[state] (ready) at (0,0) {Ready};
\node[state] (playing) at (5.6,0) {Playing};
\node[state] (paused) at (5.6,-4) {Paused};
\node[state] (completed) at (0,-4) {Completed};
\node[state] (error) at (5.6,-8) {Error};
\draw[->] (ready.east) -- node[above]{Start} (playing.west);
\draw[->] (playing.south) -- node[right]{Pause} (paused.north);
\draw[->] (paused.east) -- (8.1,-4) -- node[right]{Resume} (8.1,0) -- (playing.east);
\draw[->] (playing.north) -- (5.6,2) -- node[above]{End of recording} (-2.4,2) -- (-2.4,-4) -- (completed.west);
\draw[->] (paused.west) -- node[above,align=center,text width=2.4cm]{Step at end:\\completion policy} (completed.east);
\draw[->] (paused.south) -- node[right]{Playback failure} (error.north);
\draw[->] (playing.east) -- (9.5,0) -- node[right]{Playback failure} (9.5,-8) -- (error.east);
\node[draw,dashed,align=left,text width=4.2cm] at (-0.3,-8)
 {While paused: step admits one event. Seek restores state before continuing. Speed, seek, timer and end-state controls follow the approved policy.};
\node[align=center,text width=10cm] at (3.5,-11.3)
 {Ready follows compatibility approval. Reopening a recording returns through compatibility evaluation; detailed recovery transitions remain open.};
\end{tikzpicture}
}
\caption{Candidate playback states and stakeholder control transitions.}\label{Figure:PlaybackStates}
\end{figure}
